\scanpage{81}
On va graduer $K_{u,k}^{\bullet\bullet}$ par a) le degré extérieur venant de
$\Lambda_k^*\Psi_k$ $(\deg_{\mathrm{ext}}\Gamma^i\otimes\Lambda^j=j)$
et b) le \uline{degré total} $i+j$. Ainsi $d$ augmente le degré extérieur par 1,
et invarie le degré total \deleted{\unsure{donc}}
(Le degré complémentaire est $\deg_{\mathrm{comp}}\Gamma^i\otimes\Lambda^j=i$,
il est diminué de 1 par $d$.) Muni du degré total (ou complémentaire)
$K_{u,k}^{\bullet\bullet}$ est une Algèbre graduée sur l'Anneau gradué $k\{T\}$.

\unsure{Conséquence} (de $\unsure{\gamma}$ disons de suite)

\uline{Théorème} Soit
$K^{\bullet\bullet}(X,u,k)=\Gamma(X,K_{u,k}^{\bullet\bullet})$.
Alors on a un hom can (du fait que $K_{u,k}^{\bullet\bullet}$
\deleted{\unsure{H}} est une résolution de $k\{T\}$)
\[
\begin{gathered}
H^{**\deleted{\unsure{*}}}(X,u,k)\overset{\mathrm{dfn}}=
H^{**}(K^{\bullet\bullet}(X,u,k))\\
\longrightarrow H^{**}(X,k\{T\})\\
\simeq H^{**}(X,k)\otimes_{\mathbb Z}\mathbb Z\{T\}
\end{gathered}
\]
et ceci est un isomorphisme pour les bidegrés $(i,j)$ tels que
$\deleted{\unsure{i\leq j}}j\leq i$, tandis que
$H^{i,j}(X,u,k)=0$ (et pour cause) pour $j>i$.

\margin{\figure{\begin{tikzpicture}[x=.65cm,y=.48cm]
\draw (0,4.7)--(0,0)--(5.7,.05);
\draw (0,0) circle(.06);
\fill (1,0) circle(.07); \fill (1,1) circle(.06);
\fill (2,0) circle(.07); \fill (2,1) circle(.065); \fill (2,2) circle(.09);
\fill (3,0) circle(.09); \fill (3,1) circle(.05); \fill (3,2) circle(.09); \fill (3,3) circle(.065);
\fill (4,0) circle(.06); \fill (4,1) circle(.04); \draw (4,2) circle(.06);
\fill (4,3) circle(.065); \fill (4,4) circle(.06);
\end{tikzpicture}}}

\uline{Corollaire plus général.} Soit $f:X\longrightarrow Y$ un morphisme de topos. Alors
\[
f_*K_{u,k}^{\bullet\bullet}\longrightarrow Rf_*(k\{T\})
\simeq Rf_*(k)\otimes_{\mathbb Z}\mathbb Z\{\unsure{T}\}
\]
\margin{80}
